\mwnewpart
\renewcommand\runningtitle{Section 2. Building}
\chapter{Building}
\sectioncontents

\section{In Brief}\ixs{building}

\begin{plainwords}
Programs written in 1983 were assembled by programs of 1983. RetrOS/LT is
built the same way: Microsoft's own assembler and linker, run inside an
emulated PC on a modern computer, turn the source into the programs on
the diskettes. The parts that are new are written for NASM, a modern
assembler, which makes the same kind of program. One script does
everything, in a few minutes, and the result is the same each time.
\end{plainwords}

\begin{dialogue}
git submodule update --init
sh scripts/build.sh
python3 scripts/mkinstall.py
\end{dialogue}

\file{build/OUTPUT} then holds every program, and \file{build/install}
the two diskette images. The tools needed are DOSBox-X, NASM, mtools,
Python 3 and \typed{patch}.

\section{The Tools}\ixs{MASM}\ixs{LINK}\ixs{DOSBox-X}

Microsoft's assembler and linker are taken from the same repository: the
MS-DOS 4.0 source carries in \file{v4.0/src/TOOLS} the Macro Assembler
5.10 and the linker 3.65 it was built with, released under the same
license. The 2.0 kit has an older assembler, MASM 1.10, but it could not
build the kit reliably. MASM 5.10 builds it once two names that it
reserves are changed (\patch{0001}). \file{scripts/dosrun.sh} runs a batch
file in DOSBox-X with no window, the work directory as drive D and the
tools as drive T.

\section{Preparing the Source}\ixs{prep.py}

The published files are not all in the form the tools want.
\file{scripts/prep.py} copies \file{v2.0/source} to \file{build/work/SRC}
and cleans each text file:

\srclist{prep_clean}{Cleaning a published file.}

Some files end in an end-of-file mark followed by padding, and the
patches would not apply over it. A few were edited in WordStar, which
marks a line it has reflowed with the high bit set on the line feed; MASM
takes \code{CR 8Ah} as part of the line, and the line that follows as a
continuation of it. After cleaning, \file{prep.py} applies the series with
\typed{patch}, adds RetrOS/LT's own files from \file{src/}, and puts back the CR LF line endings and the final
\keys{Ctrl}{Z} the 1988 tools expect.

\engnote{The patches are made against the cleaned tree with LF line
endings, so that they are readable and \typed{patch} treats them as text.
The CR LF endings are added after them.}

\section{The Kernel, COMMAND.COM and RETBIO.COM}

\file{scripts/build.sh} writes one batch file, \file{BUILD.BAT}, that
assembles every module and links each program, and runs it in DOSBox-X.
The kernel is linked from twenty-one modules, \file{COMMAND.COM} from
seventeen, and \file{RETBIO.COM} from three: \file{BIOS.ASM}, then
Microsoft's \file{SYSINIT.ASM} and \file{SYSIMES.ASM}. \typed{EXE2BIN}
makes the \file{.COM} files, \file{RETBIO.COM} with a base segment of 70h,
where the startup record loads it. Link response files are written by a
few lines of Python, so that no line is longer than LINK 3.65 accepts.

SYSINIT is built as IBM's (\patch{0001}):

\srclist{sysinit_ibm}{SYSINIT as IBM built it.}

In PC DOS 2, the program-loading function, EXEC, is part of
\file{COMMAND.COM}; the kernel's own returns ``invalid function'' until
\file{COMMAND.COM} has hooked it. SYSINIT built for MS-DOS asks the kernel
to EXEC \file{COMMAND.COM}, and fails with \shown{Bad or missing Command
Interpreter}; built for IBM, it loads \file{COMMAND.COM} itself.

\section{The Parts Written for NASM}\ixs{NASM}

Everything new but \file{RETBIO.COM} and \typed{FORMAT}'s IBM part is
assembled by NASM, before the DOSBox-X run, so that the master boot record
and startup record can be built into the programs that write them:

\srclist{build_nasm}{The NASM parts.}

Each file of \file{cmd/} becomes a \file{.COM} program of the same name,
each of \file{dev/} a \file{.SYS} driver, and \file{install/INSTALL.ASM}
\file{INSTALL.COM}. \typed{FDISK} includes the first 446 bytes of the
master boot record as \file{MBRCODE.INC}; Install includes the same. The
startup record goes into the build tree as \file{BOOTSECT.INC}, which
\typed{FORMAT} and \typed{SYS} include.

\section{What Comes Out}

\begin{tabularx}{\textwidth}{@{}P{1.6in}Y@{}}
\toprule
\file{RETBIO.COM}, \file{RETDOS.COM} & the system\\
\file{COMMAND.COM} & the command processor\\
twelve programs & Microsoft's commands, \typed{FORMAT} with IBM's part\\
\file{LINK.EXE} & copied from the 2.0 release\\
ten programs, \file{ANSI.SYS} & the commands and driver written new\\
\file{INSTALL.COM} & Install\\
\file{BASICA.EXE}, \file{BASIC.COM} & BASIC (Section~8)\\
\bottomrule
\end{tabularx}

\file{scripts/build.sh} builds BASIC last, by running
\file{scripts/buildbasic.sh}, which makes its own tree and runs MASM and
LINK in DOSBox-X again (Section~8).

\file{scripts/mkinstall.py} then makes the two diskettes (Section~6),
\file{scripts/mkboot.py} a startup diskette of any of the four sizes, and
\file{scripts/mkhd.py} a fixed disk image with a DOS partition, for the
tests.
